\section{Building an assessment}\label{sec:assessments}

An assessment -- test, final, quiz, homework -- is an \file{Assessment}
document: an assembly file that sets up a cover, then imports problems from
the bank with points attached. The skeleton, modelled on the demo's
\file{_assessment/test.tex}:

\begin{manexsource}
%! views: solutions, instructor
\documentclass[12pt,thursday,automatch]{Assessment}
\assessmentsetup{
  course = {MATH 101},
  term   = {Fall 2026},
  name   = {Term Test 1},
  date   = {16 Oct},
  instructions = {
    \textbf{\sffamily Instructions:}
    \begin{itemize}[nosep]
      \item \printtotalpoints\ points.
      \item Closed book. No calculators.
      \item Show all work for full marks.
  \end{itemize}},
}
\begin{document}
\AssessmentTitle
\newpage
\importproblem[points=5]
  {../derivatives/problems/}{chain-rule.tex}
\workspace[5]
...
\end{document}
\end{manexsource}

The class renders questions as ``Q\,$N$'' hung in the \emph{left} margin with
their point marks ``($N$ pts)'' in the \emph{right} margin -- numbers left,
marks right, so standalone questions and lettered group parts drop their marks
into the same rail. None of that appears in this manual's examples (which
render with the neutral ``Problem $N$.''\ style); build the demo's
\file{test.tex} to see the real thing.

\subsection{Class options}

\begin{itemize}
  \item \texttt{[noid]} -- for documents that collect no student identity:
    homework, in-class worksheets, MC practice. Drops the NAME/ID box and the
    deep scanner top margin. The \emph{default} is the fully-furnished
    identified exam, on purpose: forgetting \texttt{[noid]} merely
    over-furnishes homework with an unused ID box, whereas the opposite
    default could ship a real exam with no ID box at all -- the safe mistake
    is the easy one to make.
  \item \texttt{[automatch]} -- for Crowdmark's automatic name matching:
    keeps the deep top margin (the QR code still lands there) but replaces
    the NAME/ID box with a blank mid-cover band (9--18\,cm down the page, per
    spec) where students write their identity for the scanner.
  \item \texttt{[tuesday]} / \texttt{[thursday]} -- a framed day box beside
    the title: TUE on the \emph{left}, THU on the \emph{right}, so two
    section piles sort at a glance by which side the box sits on, without
    reading a word.
  \item \texttt{[legalpaper]} -- legal stock; the margins stay absolute, so
    the extra three inches become writing room.
  \item Stock \file{article} options pass through (\texttt{[12pt]},
    \texttt{[twoside]}, \dots).
\end{itemize}

About that deep top margin: Crowdmark stamps a QR code into the top band of
every page at upload, so the band stays clear of content throughout, and the
cover's NAME/ID box is lifted up \emph{into} it, beside where the QR lands.

\subsection{The cover}

\cs{assessmentsetup} takes the five fields you saw above -- \env{course},
\env{term}, \env{name}, \env{date}, \env{instructions} -- every one optional;
an unset field prints nothing, brackets and all. Anything else an exam cover
might say (time allowed, formula-sheet note, coverage) is just prose in
\env{instructions}, which needs no dedicated keys. \cs{AssessmentTitle} then
emits the cover: NAME/ID box, centred title, instructions.

Note what is \emph{not} a field: the point total. Write
\cs{printtotalpoints} in the instructions (as above) and the printed total is
summed from the imported problems themselves -- it can never disagree with
the paper. (Both this total and the ``Page $N$ of $M$'' footer are two-pass
values: a first compile shows \texttt{??}, and the build's latexmk rerun
settles it.)

In a solutions build, the cover also stamps itself
{\sffamily\bfseries\color{SolutionColor}SOLUTIONS} -- the answer key
announces itself.

\subsection{Points at the import site}

Points are attached where a problem is imported, never stored in the problem
file (section~\ref{sec:problem-files}). The grammar, in increasing order of
ceremony:

\begin{manexsource}
\importproblem[5]{DIR}{FILE}          % shorthand: points=5
\importproblem[points=5]{DIR}{FILE}   % the same, spelt out
\importproblem[points=2*]{DIR}{FILE}  % announced BONUS: off the total
\importproblem[points=5, nocount]{DIR}{FILE}  % quiet: off the total
\importproblem[points={2,1*}, workspaces={4,4}]{DIR}{FILE}
                                      % per-part (a qparts file)
\importproblem[pointseach=2]{DIR}{FILE}   % broadcast to every part
\importproblem[points=5, hint]{DIR}{FILE} % force this hint ON
\end{manexsource}

\begin{itemize}
  \item \textbf{Values are integers} -- the total accumulator and the
    singular/plural logic both require it; \texttt{points=2.5} errors.
  \item \textbf{The \texttt{*} suffix means announced bonus}: the value is
    shown (a bonus \emph{is} worth $N$ if earned), labelled \emph{bonus} in
    the margin, and kept out of \cs{printtotalpoints}. In a per-part list,
    star just the bonus part: \texttt{points=\{2,3,1*\}} makes only part (c)
    bonus.
  \item \textbf{\texttt{nocount} is the quiet off-books case} -- marker
    unchanged, total untouched. Its main use is choose-$N$-of-$M$ accounting,
    where \env{problemchoice} below applies it for you.
  \item \textbf{One seam to know}: the lone-integer shorthand only fires when
    the whole option is one integer, so \texttt{[5, hint]} errors -- write
    \texttt{[points=5, hint]}.
  \item \texttt{hint} / \texttt{nohint} force one problem's hint on or off
    against the document toggle -- a worked example on a test, or a hint
    withheld in an otherwise-hinted homework.
\end{itemize}

\subsection{Grouping: \env{problemgroup} and \env{problemchoice}}

\env{problemgroup} assembles \emph{independent} one-file problems into one
numbered question with lettered parts -- the parts stay atomic files in the
bank, and each import inside carries its own points, exactly as anywhere
else:

\begin{manexsource}
\begin{problemgroup}{Evaluate}
  \importproblem[points=5, hint]
    {../derivatives/problems/}{quotient-rule.tex}
  \workspace[3]
  \importproblem[5]
    {../derivatives/problems/}{second-derivative.tex}
  \workspace[2]
\end{problemgroup}
\end{manexsource}

This renders as ``Q\,$N$.\ Evaluate'' with parts (a), (b); a \cs{ref} to a
part reads ``$N$(a)''. To \cs{label} the \emph{group}, put the label just
inside \env{problemgroup}, before the first part -- after \cs{end} it would
catch the last part instead. (Genuinely dependent sub-questions are one file
with \env{qparts}, section~\ref{sec:problem-files} -- import that file here
like any other.)

\env{problemchoice} is the choose-$N$-of-$M$ wrapper:

\begin{manexsource}
\begin{problemchoice}[select=2, points=5]
  \importproblem[]
    {../derivatives/problems/}{mc-domain.tex}
  \FreshPage
  \importproblem[]
    {../derivatives/problems/}{mc-tangent-slope.tex}
  \FreshPage
  \importproblem[]
    {../derivatives/problems/}{mc-select-all.tex}
\end{problemchoice}
\end{manexsource}

Each member question repeats the generated instruction (``Answer any two of
Questions $N$--$M$.''\ -- page-break safe), every member shows the shared
per-question value, and the total counts it $\text{select} \times
\text{points}$ once, not per member. Members therefore carry no
\texttt{points=} of their own (one that does triggers a warning, since it
would silently diverge from the shared value).

\texttt{instruction=} says \emph{where} that sentence appears, and the
default is the block above: \texttt{each}. Repeating it earns its keep when
the choice is one block among required questions -- a student meeting Q8
alone still learns that it is one of a choice. It is noise when the
\emph{whole paper} is one choice, where the sentence is really a fact about
the paper: then set \texttt{instruction=first} to announce the run once, or
\texttt{instruction=none} and say it in the cover instructions.

\begin{manexsource}
\begin{problemchoice}[select=5, points=6, instruction=none]
\end{manexsource}

\noindent
Nothing checks that the paper says it somewhere -- this package cannot read
your cover. Both \cs{label}s are still dropped in every mode, so a cover line
can name the range itself with \cs{ref}\texttt{\{pc@1@first\}} and
\cs{ref}\texttt{\{pc@1@last\}} (the digit is the group's number, counting
from one).

\subsection{Answer space and page furniture}

\cs{workspace}\texttt{[N]} reserves $N$ lines of handwriting after a question
(also \texttt{[halfpage]}, \texttt{[fullpage]}, or \cs{workspace*} to fill
the rest of the page); \cs{FreshPage} is a collapsible page break;
\cs{FreshPage}\texttt{[blank]} and \texttt{[watermark]} emit a dedicated
blank page, the latter stamped ``Blank page for extra work.'' All of this is
Workspace machinery with one deep design idea -- reserved space that
solutions can paint into without moving a single page break -- explained in
section~\ref{sec:overlay}.

The class's own page furniture is minimal and load-bearing: a ``Page $N$ of
$M$'' footer on every page (a student can check the stapled packet is
complete at a glance) and no running head.

\subsection{Views of one assessment}

With \texttt{\%! views: solutions, instructor} in the magic comment, the
demo's \file{test.tex} builds three PDFs: the student paper, an answer key
(\file{-SOLUTIONS}), and the staff bundle (\file{-INSTRUCTOR}: solutions +
hints + instructor notes, inside the page-edge \emph{Instructor copy} frame).
Under the default overlay mode all three have \emph{identical page breaks} --
the whole point of section~\ref{sec:overlay}. (Add \texttt{hints} to the list
-- as the demo's \file{quiz.tex} does -- for a fourth, student-facing
\file{-HINTS} copy.) Build mechanics are in section~\ref{sec:build-and-find}.

\subsection{Printing the cover on its own}

A test's first page is the one page you may want to print by itself --- on
heavier paper, or in a different quantity from the paper behind it. Adding
\texttt{\%! coverpage: yes} to the magic comment writes
\file{test_COVERPAGE.pdf} beside \file{test.pdf}: page~1 of the finished
test, nothing else. It is taken from the student build only, since a
solutions or instructor cover is the same sheet of paper.

Note \emph{finished}. The obvious way to do this would be a page-selection
package in the preamble, and it does not work here, for a reason worth
knowing because it applies to anything that stops a run early. The cover's
own numbers are two-pass values: \cs{printtotalpoints} is a
\texttt{totcount} sum written to the \file{.aux} at \cs{end\{document\}} and
read back on the next run, and the footer's ``Page 1 of $M$'' is
\cs{pageref\{LastPage\}}. The package for the job, \texttt{pagesel},
suppresses the \file{.aux} outright and aborts the job after the last
selected page --- so the run never reaches the point where the total is
written, and the cover prints the \emph{wrong} number of points. Cutting
page~1 out of a normally-built PDF has no such failure mode: the document
compiled exactly as it always does, and the page is the page.
